% quickstart 玩具样例 · 极简中文论文（桩）
% 编译：xelatex -interaction=nonstopmode main.tex
% 玩具题：产品A 单位售价 p=50，成本 C(x)=a x^2-10x+100（标称 a=1），
%         求利润最大产量，并在 a 取 0.9/1.0/1.1（±10%）下做灵敏度分析。
\documentclass[11pt]{article}
\usepackage[UTF8]{ctex}
\usepackage{amsmath}
\usepackage{unicode-math}
\usepackage{booktabs}
\usepackage{tabularx}
\usepackage{geometry}
\usepackage{needspace}
\geometry{a4paper, margin=2.5cm}
\IfFontExistsTF{Times New Roman}{%
  \setmainfont{Times New Roman}%
}{%
  \setmainfont{TeX Gyre Termes}%
}
\setmathfont{TeX Gyre Termes Math}

\title{产品 A 单变量生产优化（quickstart 玩具样例）}
\author{工具链自检}
\date{\today}

\begin{document}
\maketitle

\section{摘要}
针对问题1，本文以产品 A 的生产利润最大化为目标，建立单变量二次规划模型。
由售价与二次成本构造利润函数后，利用严格凹性确定唯一最优产量；
成本二次项系数再分别取标称值上下浮动百分之十，按同一解析式重算最优产量与最大利润，
用以评价方案的鲁棒性。结果表明最优产量随成本系数增大而平稳下降，
标称情形得到\textbf{最优产量 $30$ 单位、最大利润 $800$ 元}；在成本系数正负百分之十扰动下，
\textbf{最优产量保持在 $[27.27,\,33.33]$ 单位、最大利润保持在 $[718.18,\,900.00]$ 元}，
各情景的最优产量均满足非负约束；区间是否满足实际稳健性要求，仍需结合决策者可接受的偏差阈值判断。

\noindent 关键词：单变量优化；二次规划；灵敏度分析；鲁棒性

\section{引言}
\textbf{问题1}：某工厂生产产品 A，单位售价固定为 50 元，成本随产量呈二次形式变化；
在成本二次项系数存在正负百分之十扰动的条件下，求使利润最大的最优产量，并评价该方案的鲁棒性。

\section{问题分析}
本题是带非负域约束的单变量凹二次优化：输入为售价 $p=50$、成本系数
$a\in\{0.9,1.0,1.1\}$ 和二次成本形式，决策变量为满足 $x\geq0$ 的产量，目标是利润最大化，
输出为各情景的 $x^*$ 与 $P^*$。给定参数范围内的解析驻点均为正，边界约束不激活；先求唯一最大值，
再对 $a$ 做同口径灵敏度分析。每个情景给出唯一可行最优点及可复核区间即完成求解；
没有业务阈值时不直接宣称“足够鲁棒”。本示例仅有一问，不存在跨问的数据、参数或误差传递。
解析优化与灵敏度分析的基本建模口径参考文献\cite{a,b}。

\section{模型假设}
为简化问题，本文做出以下假设：

假设1：单位售价固定不变，产量与销量相等，不考虑库存积压、产能上限与设备折旧。
这一假设用于隔离成本系数扰动对最优解的影响，与题面给出的固定售价口径一致；若现实中存在滞销或产能瓶颈，
本文得到的最优产量和利润将偏高，因此结论只适用于上述简化边界。

\Needspace{12\baselineskip}
\section{主要符号说明}
本文使用的主要符号如表~\ref{tab:sym} 所示。

\begin{table}[htbp]
  \centering
  \caption{主要符号说明}
  \label{tab:sym}
  \begin{tabularx}{\textwidth}{>{\centering\arraybackslash}p{0.16\textwidth}X>{\centering\arraybackslash}p{0.22\textwidth}}
    \toprule
    符号 & 含义或定义 & 单位/取值域 \\
    \midrule
    $x$ & 产品 A 的产量 & 单位 \\
    $p$ & 产品 A 的单位售价 & 元/单位 \\
    $a$ & 成本二次项系数，标称值为 $1.0$ & 元/单位$^2$ \\
    $P(x)$ & 产量为 $x$ 时的利润 & 元 \\
    \bottomrule
  \end{tabularx}
\end{table}

\section{问题一：最优产量与鲁棒性}
本问继承题面给定的售价、二次成本形式与模型假设，需要同时回答标称最优产量和成本扰动下的稳定性。
输入为 $p=50$、$a\in\{0.9,1.0,1.1\}$，产量受 $x\geq0$ 约束；先解析求取利润函数的唯一最大值，
再对成本系数做正负百分之十的同口径重算。最终给出各情景的最优产量、最大利润及其扰动区间；
所有最优点均可行且计算可复核，即满足本问的求解标准。本问是唯一问题，不向后续问题传递结果。

\subsection{模型建立}
设产量为 $x\geq 0$，单位售价 $p=50$，成本函数 $C(x)=a x^2-10x+100$（标称 $a=1$）。
利润为销售收入减去成本：
\[
P(x)=p\,x-C(x)=50x-(a x^2-10x+100)=-a x^2+60x-100.
\]
由于 $a>0$，二次项系数为负，$P(x)$ 是开口向下的严格凹二次函数，唯一驻点即全局最大值点。

\subsection{模型求解}
令 $P'(x)=-2a x+60=0$，得最优产量 $x^{*}=30/a$。代入标称值 $a=1$，
$x^{*}=30$，最大利润 $P(x^{*})=-30^{2}+60\cdot30-100=800$ 元。

\subsection{模型检验与分析}
优化模型的最优解依赖于成本二次项系数 $a$ 的取值，而实际生产中原材料价格、能耗与人工成本
都会使 $a$ 在一定范围内波动，若不考察扰动下的解稳定性，所给出的最优产量在真实环境中可能失效，
因此有必要对 $a$ 做灵敏度分析。

\subsubsection{灵敏度分析}
将成本二次项系数 $a$ 分别取 0.90、1.00、1.10，对应标称值上下浮动百分之十，
重新代入 $x^{*}=30/a$ 与 $P^{*}=900/a-100$ 求解，结果如表\ref{tab:sen}所示。

\begin{table}[htbp]
  \centering
  \caption{成本系数扰动下的灵敏度分析结果}
  \label{tab:sen}
  \begin{tabular}{ccc}
    \toprule
    成本系数 $a$ & 最优产量 $x^{*}$ & 最大利润 $P^{*}$（元） \\
    \midrule
    0.90 & 33.33 & 900.00 \\
    1.00 & 30.00 & 800.00 \\
    1.10 & 27.27 & 718.18 \\
    \bottomrule
  \end{tabular}
\end{table}

由表\ref{tab:sen}可见，当 $a$ 在标称值上下浮动百分之十时，
最优产量在 $[27.27,\,33.33]$ 区间内移动，最大利润在 $[718.18,\,900.00]$ 元之间变化。
相对标称利润 800 元，成本系数下降百分之十时利润上升约百分之十二点五，
成本系数上升百分之十时利润下降约百分之十点二。上述区间给出了方案对该参数扰动的实际响应幅度；
其稳健性是否可接受，应由生产决策中预先设定的产量与利润偏差阈值判定。

\section{模型的评价}
优点1：模型结构简单，最优解可解析求得，无数值迭代误差。

优点2：通过灵敏度分析定量给出了最优方案在参数扰动下的变化区间，便于按业务阈值判断稳健性。

缺点1：单变量与二次成本假设会遗漏多产品耦合和产能瓶颈，可能低估高负荷区间的边际成本并高估
可实现利润；实际应用时应加入产能约束、交叉成本项并重新检验最优解。

\section{模型的改进}
后续可将模型推广为多产品、多约束的线性或非线性规划，并采用数值优化算法求解，
同时对更多成本参数做联合扰动分析。

\section{模型的推广}
原问题只有标量产量、非负约束和三个确定性成本情景，评价量是最优产量与最大利润。
推广到多产品生产时，决策变量变为产量向量，并需加入产能、库存和产品耦合约束；价格与成本也应由历史样本形成的
经验分布或压力情景替代三个固定取值。相应评价指标除期望利润外，还应包含约束违反率、利润下行风险及解的稳定区间。
因此迁移时必须重新标定参数、重建约束，并用新数据重新验证可行性和敏感性，不能直接照搬本例的数值结论。

\clearpage
\begin{thebibliography}{9}
\bibitem{a} 姜启源, 谢金星, 叶俊. 数学模型[M]. 5版. 北京: 高等教育出版社, 2018.
\bibitem{b} 司守奎, 孙玺菁. 数学建模算法与应用[M]. 2版. 北京: 国防工业出版社, 2021.
\end{thebibliography}

\end{document}
